\section{Restriction after every prime scalar is killed}
\label{sec:all-cm-original-restriction}

Define
\[
 R_{\rm all}=R_{\rm quad}\cap\bigcap_{p\geq5}\ker h_p,\qquad
 V_{\rm all}=V_{\rm quad}\cap R_{\rm all}.
\]
For any closed normal subgroup \(L\subseteq N\), write
\[
 B(L)=L/\overline{[G,L](L\cap H)}.
\]
These are intrinsic relative quotients of the indicated
groups. Their comparison with \(B\) requires attention
to which characters extend.

\begin{theorem}
\label{thm:all-cm-original-restriction}
The original quadratic restriction maps satisfy
\begin{equation}
 \ker\bigl(X_2\to H^1(R_{\rm all},\mathbb F_2)^G\bigr)
 =\ker\bigl(X_2\to H^1(V_{\rm all},\mathbb F_2)^G\bigr)
 =\mathcal S.
 \label{eq:all-cm-original-kernel}
\end{equation}
There are canonical maps
\begin{equation}
 B(V_{\rm all})_{\max2}
       \xrightarrow{\ \sim\ }B(R_{\rm all})_{\max2}
       \twoheadrightarrow B_{\max2}.
 \label{eq:all-cm-relative-surjection}
\end{equation}
The last map is a surjection; an isomorphism is not
asserted. Moreover,
\[
 R_{\rm all}/V_{\rm all}
       \simeq R_{\rm quad}/V_{\rm quad},
\]
and this isomorphism preserves the image of the
arithmetic intersections.
\end{theorem}
\begin{proof}
Let \(C\) be the image of the joint scalar homomorphism
\[
 h:R_{\rm quad}\longrightarrow\prod_{p\geq5}P_p.
\]
It is a compact abelian pro-\(2\) group.
The isolated local arithmetic axes and compactness give
\begin{equation}
 A_0:=\prod_{p\geq5}Q_pP_p
                   \subseteq h(V_{\rm quad}\cap H)\subseteq C.
 \label{eq:all-cm-arithmetic-axes}
\end{equation}
The image of \(V_{\rm quad}\) equals \(C\): the quotient
of these two images is both a quotient of the pro-odd
group \(R_{\rm quad}/V_{\rm quad}\) and a pro-\(2\)
group, so it is trivial.
The same argument applied to
\((R_{\rm quad}\cap H)/(V_{\rm quad}\cap H)\)
shows that their arithmetic images under \(h\) also agree.

We claim
\begin{equation}
 C=A_0+\overline{[G,C]}.
 \label{eq:all-cm-projection-identity}
\end{equation}
Here additive notation is used in the compact scalar group.
For each \(p\), let \(S_p\) be conjugation by a primitive
arithmetic \(p\)-inertia element. Equivariance makes
\(S_p\) act as \(s_p\) on the \(p\)-coordinate and as
the identity on every other prime coordinate.
The reflex polynomial has augmentation
\(\epsilon_p=Q_p(1)\in\{1,-1\}\).
Write
\[
 Q_p(s)=\epsilon_p+(s-1)F_p(s),\qquad F_p\in\mathbb Z[s].
\]
For \(c=(c_q)_q\in C\), let \(i_p\) denote the
single-coordinate inclusion into the ambient product.
There is an identity in that product:
\begin{equation}
 \epsilon_p i_p(c_p)
   =i_p(Q_pc_p)-(S_p-1)F_p(S_p)c.
 \label{eq:all-cm-coordinate-projection}
\end{equation}
The first term on the right belongs to \(A_0\),
and the second to \([G,C]\).
Thus \(i_p(c_p)\) belongs to the right side of
\eqref{eq:all-cm-projection-identity}, as well as to \(C\).
Finite sums of these coordinate projections converge
to \(c\). The right side of
\eqref{eq:all-cm-projection-identity} is closed, being
the sum of two compact subgroups. This proves the claim.
In particular every continuous \(G\)-invariant
character of \(C\) which vanishes on its arithmetic
image is zero.

Use the established decomposition
\(X_2=\operatorname{Hom}(B,\mathbb F_2)\oplus\mathcal S\).
The summand \(\mathcal S\) vanishes on \(R_{\rm quad}\).
If a relative character from \(\operatorname{Hom}(B,\mathbb F_2)\)
vanishes on \(R_{\rm all}\), its restriction to
\(R_{\rm quad}\) factors through \(C\), is invariant,
and kills the arithmetic image containing \(A_0\).
It is zero by \eqref{eq:all-cm-projection-identity};
the known factorization \(G=R_{\rm quad}H\) makes the
original character zero. The same proof on \(V_{\rm quad}\)
works because its joint image is \(C\) and its original
relative-character restriction is an isomorphism.
This proves \eqref{eq:all-cm-original-kernel}.

Equality of the full joint images gives
\(R_{\rm quad}=R_{\rm all}V_{\rm quad}\), hence
\[
 R_{\rm all}/V_{\rm all}\simeq R_{\rm quad}/V_{\rm quad}.
\]
Equality of the joint arithmetic images allows the same
matching using \(V_{\rm quad}\cap H\), proving the
assertion about arithmetic images in this quotient.
It is pro-odd. Coprime extension of invariant
characters with finite abelian \(2\)-group values,
including their zero arithmetic restriction, therefore
gives the first isomorphism in
\eqref{eq:all-cm-relative-surjection}.

Finally, the canonical image of
\(B(R_{\rm all})_{\max2}\) in \(B_{\max2}\) is closed.
If it were proper, the resulting nonzero abelian
pro-\(2\) quotient would have a nonzero continuous
character to \(\mathbb F_2\). Its pullback would be
an original relative quadratic character vanishing
on \(R_{\rm all}\), contrary to the injectivity just
proved. The last arrow is consequently surjective.
\end{proof}

For clarity, the augmentation used above is integral.
Represent the inverse CM type as
\[
 \Phi_p^*=(1+s+\cdots+s^{h-1})+(c-1)A(s),
 \qquad\deg A<h,
\]
where each coefficient of \(A\) is \(0\) or \(1\).
The identity
\((s-1)\Phi_p^*=(c-1)(1+(s-1)A(s))\)
gives \(Q_p\) after reduction by \(s^h+1\).
Its augmentation is \(1-2a_{h-1}\), hence is \(1\)
or \(-1\). This also explains why no determinant
or irreducibility hypothesis is needed in
\eqref{eq:all-cm-coordinate-projection}.

\subsection{The intrinsic extension problem}

Let \(R_0=R_{\rm quad}\), let \(C=h(R_0)\), and put
\(D_0=h(R_0\cap H)\). Thus \(R_{\rm all}=\ker h\),
\(G=R_0H\), and
\[
 C=D_0+\overline{[G,C]}
\]
by \eqref{eq:all-cm-projection-identity}.
The following criterion separates the central extension class
from the conjugation action on its possible splittings.
Theorem~\ref{thm:all-cm-central-extension} will eliminate
the second condition for this CM quotient.

\begin{proposition}
\label{prop:all-cm-intrinsic-extension}
Let \(\kappa:R_{\rm all}\to\mathbb F_2\) be continuous,
\(G\)-invariant, and zero on \(R_{\rm all}\cap H\).
There is a central relative obstruction
\[
 o_2(\kappa)\in H^2(C,D_0;\mathbb F_2).
\]
If it vanishes, there is a conjugation obstruction
\[
 o_1(\kappa)\in
 H^1_{\rm cont}\!\left(G/R_0,\
    \operatorname{Hom}_{\rm cont}(C/D_0,\mathbb F_2)\right).
\]
The character extends to a continuous character of
\(G\) which is zero on \(H\) if and only if both
obstructions vanish. When it exists, the extension
is unique. Characters on \(V_{\rm all}\) satisfy
the same criterion after their unique coprime
extension to \(R_{\rm all}\).
\end{proposition}
\begin{proof}
For nonzero \(\kappa\), form the central extension
\(R_0/\ker\kappa\to C\) with kernel \(\mathbb F_2\).
Its kernel is central because \(\kappa\) is
\(G\)-invariant. The image of \(R_0\cap H\) gives
a prescribed continuous splitting over \(D_0\);
its kernel is \(R_{\rm all}\cap H\), on which
\(\kappa\) is zero. Define \(o_2(\kappa)\) to be
the relative extension class, and set \(o_2(0)=0\).
This class vanishes exactly when the extension has
a continuous splitting agreeing with the prescribed
one. Projection onto the central factor is then
a character of \(R_0\) extending \(\kappa\) and
zero on \(R_0\cap H\). Conversely the kernel of
such a character gives the required splitting.

These relative classes use continuous cochains
without an extra topological hypothesis.
A finite-kernel profinite extension has a continuous
section: choose an open subgroup meeting the kernel
trivially and choose representatives for its finitely
many image cosets. The section can be made to agree
with a prescribed section over a closed subgroup,
because a continuous finite-valued function on a
closed subset of a profinite space extends by
clopen separation. This realizes the extension
and its splitting in continuous relative cohomology.

When \(o_2(\kappa)=0\), the possible extensions to
\(R_0\) form a torsor under
\(\mathscr U=\operatorname{Hom}_{\rm cont}(C/D_0,\mathbb F_2)\).
Conjugation by \(G\) preserves this torsor: the subgroup
\(H\) preserves \(R_0\cap H\), and inner conjugation
by \(R_0\) acts trivially on characters of \(R_0\).
The action factors through \(G/R_0\).

Give these character spaces the discrete topology.
Every continuous character \(\lambda:R_0\to\mathbb F_2\)
has open stabilizer: choose an open normal
\(W\triangleleft G\) with
\(W\cap R_0\subseteq\ker\lambda\); then
\([W,R_0]\subseteq W\cap R_0\), so \(W\) fixes
\(\lambda\). This proves continuity even when
\(\mathscr U\) is infinite.
Choosing one extension defines the continuous
affine cocycle of the torsor, whose cohomology
class is \(o_1(\kappa)\). It vanishes exactly when
there is a \(G\)-invariant extension to \(R_0\).

Finally a \(G\)-invariant character of \(R_0\)
zero on \(R_0\cap H\) extends along \(G=R_0H\)
by \(\chi(rh)=\chi(r)\), continuously by compact
descent. The difference of two possible invariant
extensions belongs to \(\mathscr U^{G/R_0}\).
This fixed space is zero because
\(C=D_0+\overline{[G,C]}\).
Thus the extension is unique.

The pro-odd quotient
\(R_{\rm all}/V_{\rm all}\) gives the asserted
unique extension from \(V_{\rm all}\), including
the zero arithmetic restriction, by continuous
coprime cohomology.
\end{proof}

For an isomorphism of the relative pro-\(2\) groups
in \eqref{eq:all-cm-relative-surjection}, the extension
property must hold for invariant relative characters
with every finite abelian \(2\)-group as target.
The proposition gives the exact first case, with
target \(\mathbb F_2\). Theorem~\ref{thm:all-cm-central-extension}
proves that its conjugation obstruction always vanishes.
The relative central class remains to be determined;
the factorization \(G=R_{\rm all}H\) is not asserted.
